\documentclass[11pt]{article}
\usepackage[margin=1.05in]{geometry}
\usepackage{amsmath, amssymb}
\usepackage{booktabs}
\usepackage{multirow}
\usepackage{xcolor}
\usepackage{hyperref}
\hypersetup{colorlinks=true, linkcolor=blue!50!black, urlcolor=blue!50!black}

\newcommand{\code}[1]{\texttt{#1}}
\DeclareMathOperator{\hn}{HN}

% Table bodies are generated by
% code/analysis/time_block_forecast/simstudy/expA_breakdown.R and are
% never edited by hand; this report transcribes no number.
\newcommand{\tblpath}{../../code/analysis/time_block_forecast/simstudy/output/tables}
\newcommand{\tbl}[1]{\input{\tblpath/#1_body.tex}}

\title{\textbf{Full-series threshold Brier scores for the time-block
forecast}\\[2pt]
\large Per-block and per-dataset breakdowns, Experiment~A, settings 4, 5
and 7}
\author{Score report}
\date{10 August 2026}

\begin{document}
\maketitle

\begin{abstract}
The time-block Brier score is now computed against thresholds drawn from
each dataset's \emph{full series}, replacing the validation block's own
quantile. This report gives the breakdown that a pooled lead-time table
hides: the same scores resolved by block and by dataset, at two
co-primary thresholds, for the three settings fitted under this campaign.
Those three span persistence end to end --- a short-memory control
(setting~4), a near-unit-root process (setting~5) and a long-memory
process (setting~7) --- so the score can be read against the property it
is supposed to be sensitive to. At the near-unit-root setting~5 the AR(2)
advantage is not an aggregate artifact --- AR(2) scores below i.i.d.\ in
\textbf{all five blocks} and, at $q_{90}$, in \textbf{all ten datasets},
and the pooled contrast is significant at every threshold from $q_{90}$
to $q_{98}$ (AR(2)$-$AR(1) at $q_{90}$: $t\,p = 0.0002$). This reverses
the conclusion of the retired campaign, which held that the Brier score
was too noisy to detect what CRPS could see; under the full-series rule
the Brier score alone resolves setting~5 decisively, which is why this
report carries no CRPS. The control gives the other end: at setting~4,
whose autocorrelation decays in six days, the pooled AR(2)$-$i.i.d.\ gap
over the fifteen-day horizon is $-1.0 \times 10^{-6}$ at $q_{90}$ and
stays below $0.08 \times 10^{-3}$ in absolute value at every one of the
eleven thresholds. The score therefore moves with persistence, not with
the parameter count. At the long-memory setting~7 the score is a genuine
null, and \S\ref{sec:tie} shows the null is \emph{not} an artifact of
threshold degeneracy: it survives at $q_{90}$, where the share of cells
with no realised exceedance falls from $60.9\%$ to $41.6\%$. What does
move at setting~7 is calibration --- the AR(2) predictive over-states
the exceedance probability by a factor rising to four in the far tail,
which is a level error the pooled score never shows.
\end{abstract}

\section{Scope and design}

Every number here comes from the three score caches
\code{simstudy/output/results/scores\{4,5,7\}\_hn.RData}, which hold the
Experiment~A campaign: expanding-window blocks with seams at
$T_o = 50, 80, 110, 140, 170$, each forecasting $H = 15$ days; three
methods (i.i.d., AR(1), AR(2)); ten datasets; five blocks. All fits carry
the $\lambda \sim \hn(0, 20)$ prior, and the caches are complete --- no
cell is missing, so no clean-$n$ selection is applied anywhere in this
report and $n = 10$ throughout. The three settings are fitted under one
protocol, so nothing but the generating persistence separates them.

\begin{center}
\small
\begin{tabular}{clll}
\toprule
id & label & temporal structure & persistence \\
\midrule
4 & control & AR(2) $\phi = (0.80, -0.35)$ & $\rho(F) = 0.592$,
$h^\ast = 6$ d (short memory) \\
5 & near-unit-root & AR(2) $\phi = (0.15, 0.80)$ & $\rho(F) = 0.973$,
$h^\ast = 109$ d (short memory) \\
7 & LM $d = 0.45$ & ARFIMA$(0, 0.45, 0)$ & Hurst $0.95$, hyperbolic ACF
(long memory) \\
\bottomrule
\end{tabular}
\end{center}

\noindent Setting~4 is the design's negative control. Its memory horizon
is $h^\ast = 6$ days, well inside the $15$-day forecast window, so by
lead $7$ the AR(2) and i.i.d.\ predictives have converged and the
comparison has nothing left to detect. A campaign that reported a gain
here would be reporting an artifact of the extra parameters; the value of
the setting lies in the null it is expected to produce.

\paragraph{Two co-primary thresholds.} Tables report $q_{90}$ and
$q_{95}$ side by side rather than $q_{95}$ alone. They are not
interchangeable: the two disagree on how many datasets favour AR(2) at
every setting ($10/10$ vs.\ $9/10$ at setting~5; $6/10$ vs.\ $4/10$ at
setting~7; $5/10$ vs.\ $4/10$ at setting~4), and $q_{90}$ has materially
better coverage, which \S\ref{sec:tie} needs. The full eleven-point grid
is in the appendix.

\paragraph{What is not here.} Setting~6 (ARFIMA $d = 0.35$) and
Experiments~B and~C are absent. They were fitted under the default
$\lambda \sim N(0, 20)$ prior, which is a different arm of the campaign,
and adding them would mean mixing priors inside one table rather than
re-scoring an existing cache. Their fits do survive on the machine that
produced them --- \code{simstudy/results/} holds settings 1--4 and
\code{block1\_positive\_control/results/} holds the $160$ block-1 fits of
the guarded campaign --- but every one of them carries the wrong prior for
this report, and the block-1 fits carry a different protocol as well
(one block, a fit-level gate, and a clean-$n$ selection). Those
experiments, and the CRPS results for all of them, remain recorded in
\code{tex/timeblock\_expABC\_legacy}, which is frozen for that reason.
Setting~4 was in this list until 2026-08-09, when it was refitted from
scratch under the present protocol; it is now a full member of the
campaign.

\paragraph{Brier only.} This report and its sibling
\code{tex/timeblock\_brier\_blockq} are Brier-only by construction. CRPS
is invariant to the threshold rule, so it cannot discriminate between the
two documents; and, as the abstract notes, it is no longer needed to
carry the setting-5 result.

\section{The full-series threshold rule}

For dataset $d$ at setting $s$ the thresholds are the quantiles of the
whole record,
\[
  \tau_p(d,s) \;=\; \mathrm{quantile}\bigl(y[\,\cdot,\cdot,d,s\,],\,p\bigr),
  \qquad p \in \{0.90,\ldots,0.995\},
\]
fixed across every block, lead and method, and identical for the methods
being compared. This is the rule the spatial simstudy already used
(\code{code/analysis/simstudy/scores.R}), so the two studies are now
commensurable.

The rule it replaced took the threshold from the validation window
itself, $\mathrm{quantile}(y_{\mathrm{val}}, p)$. That is look-ahead --- the
threshold is a statistic of the very window being forecast --- and it has
a second, sharper defect quantified in the sibling report: it pins the
realised exceedance rate at exactly $1-p$ in every block, so no level
error can move the base rate. The scorer still computes it, as
\code{brier.lead.blockq}, and the caches record which rule produced which
array in \code{brier\_threshold\_basis}; the merge gate refuses to
combine caches written under different rules, so an old-rule chunk cannot
silently average into a new-rule campaign. \code{expA\_breakdown.R}
re-checks that tripwire on load.

\section{Per-block breakdown}
\label{sec:block}

\tbl{expA_block_full}

\noindent\textbf{Setting 5.} AR(2) scores below i.i.d.\ in all five
blocks at both thresholds. The advantage is therefore not an artifact of
pooling: no single window carries it. What the breakdown also exposes is
how unequal the windows are --- at $q_{95}$ the i.i.d.\ level ranges from
$0.0154$ (block 3) to $0.1060$ (block 5), nearly sevenfold, and at
$q_{90}$ from $0.0451$ to $0.1759$. That spread is the rule working as
intended. A fixed threshold lets a quiet window be quiet and a turbulent
one be turbulent; the block-quantile rule flattens exactly this, which is
what the sibling report is about. The largest gap sits in the most
turbulent window (block 5, $-55.4 \times 10^{-3}$ at $q_{90}$), i.e.\ the
AR(2) memory pays where there is something to predict.

\noindent\textbf{Setting 7.} Mixed in sign --- AR(2) ahead in three of
five blocks at both thresholds, behind in block 2 by more than it leads
anywhere. No block-level structure supports an advantage.

\noindent\textbf{Setting 4.} AR(2) is ahead in one block of five at both
thresholds (block 3, $-5.6 \times 10^{-3}$ at $q_{90}$) and marginally
behind in the other four, none of them by more than $3.6 \times 10^{-3}$.
The pooled gap is $-1.0 \times 10^{-6}$ because one block's lead cancels
four small deficits, which is what scatter around zero looks like when it
is resolved by block. The contrast with setting~5 is the point of putting
the two in the same table: identical protocol, identical estimator, and
five aligned signs there against a cancelling five here.

\section{Per-dataset breakdown}
\label{sec:dataset}

\tbl{expA_dataset_full}

\noindent The dataset is the pairing unit of the tests in
\S\ref{sec:tests}, so this table shows what those $p$-values are made of.
At setting~5 and $q_{90}$, AR(2) wins in \emph{all ten} datasets, with
gaps from $-0.0787$ down to $-0.0001$; at $q_{95}$ it wins in nine, the
exception being dataset~10 at $+0.0008$. Dataset~10 is in fact the least
favourable dataset in all four (setting, threshold) combinations, which
is worth recording as a single consistent outlier rather than noise.
Neither result depends on one or two datasets carrying the mean.

At setting~7 the count is $6/10$ at $q_{90}$ and $4/10$ at $q_{95}$ ---
indistinguishable from a coin, and the individual gaps are an order of
magnitude smaller than setting~5's.

Setting~4 reads the same way, $5/10$ and $4/10$, with the per-dataset
gaps spread symmetrically about zero (at $q_{90}$, from $-0.0078$ on
dataset~8 to $+0.0047$ on dataset~1). No dataset is a consistent outlier
across the two thresholds in the way setting~5's dataset~10 is.

\section{Paired tests}
\label{sec:tests}

\tbl{expA_tests_full}

\noindent Tests are one-sided at the dataset unit, $n = 10$: within a
dataset the lead window is averaged first (the window is the endpoint,
not a replicate axis), then the five blocks, leaving one number per
dataset. This is the unit \code{expA\_threeway.R} calls primary, and
\code{expA\_breakdown.R} reproduces that script's $q_{95}$ output to
$5 \times 10^{-16}$ as a check on the aggregation.

At setting~5 every AR(2) contrast is significant at both thresholds and
both lead windows, on both the $t$ and Wilcoxon tests; the sharpest is
AR(2)$-$AR(1) over leads $1$--$15$ at $q_{90}$, $t\,p = 0.0002$. That the
\emph{second-lag} contrast is the sharpest is the expected signature of
$\phi = (0.15, 0.80)$, where a single lag cannot represent the process.
AR(1)$-$i.i.d.\ is correspondingly the only contrast that does not hold
up over the full window: it separates over leads $1$--$5$ (Wilcoxon
$p = 0.007$ at $q_{90}$, $0.019$ at $q_{95}$) but not over $1$--$15$,
where it decays to $p = 0.16$ and $0.42$. One lag buys something at short
lead and nothing thereafter; two lags buy throughout.

At setting~7 nothing reaches significance at either threshold, in either
direction.

Setting~4 likewise separates nowhere, and it is worth recording how far
from separation it is: over leads $1$--$15$ the AR(2)$-$i.i.d.\ gap is
$-1.0 \times 10^{-6}$ at $q_{90}$ ($t\,p = 0.50$) and
$-7.7 \times 10^{-5}$ at $q_{95}$ ($t\,p = 0.45$), and the closest any
contrast comes to significance is AR(2)$-$AR(1) over $1$--$15$ at
$q_{95}$, at $t\,p = 0.18$. The short-lead window is the only place the
control shows even a direction --- AR(2)$-$i.i.d.\ is
$-0.94 \times 10^{-3}$ over leads $1$--$5$ at $q_{90}$ against
$-0.001 \times 10^{-3}$ over $1$--$15$ --- which is the memory horizon
showing through: what little the second lag buys is spent by lead $6$,
exactly where $h^\ast$ says it should be.

\section{Coverage and degeneracy}
\label{sec:coverage}

The cost of a fixed threshold is that a quiet window may contain no
exceedance at all. Where that happens the Brier score has no event to
discriminate and collapses onto the predicted probability alone --- it
still scores calibration, but it can no longer reward resolution. The
ledger below is computed from the data and the thresholds only, so it is
complete regardless of which fits exist.

\tbl{expA_coverage}

\noindent At $q_{95}$, $80.8\%$ of (lead, dataset, block) cells at
setting~4, $75.3\%$ at setting~5 and $60.9\%$ at setting~7 realise no
exceedance; by $p = 0.995$ this is above $92\%$ at all three. This is the
reason $q_{90}$ is carried as a co-primary threshold rather than
relegated: at $q_{90}$ the degenerate shares fall to $63.3\%$, $59.3\%$
and $41.6\%$. The ordering is worth noting --- the control has the
\emph{poorest} coverage of the three, so its null is also the least
powerful of the three, and \S\ref{sec:tie} treats it accordingly.

The predicted-rate columns are a level-error ledger. At setting~5 both
methods under-predict exceedance relative to what is realised
($0.1189$ and $0.1027$ against $0.1257$ at $q_{90}$), AR(2) slightly more
so. Setting~7 inverts, and this is the more interesting reading: AR(2)
\emph{over}-predicts throughout ($0.1250$ against a realised $0.1111$ at
$q_{90}$), and the discrepancy widens with the threshold until at
$p = 0.995$ it predicts $0.0234$ against a realised $0.0057$ --- a
factor of four. Under long memory the fitted AR(2) puts too much mass in
the upper tail. A pooled Brier score cannot show this; the ledger can.
Setting~4 sits between the two and is mild in both directions: both
methods over-predict slightly at $q_{90}$ ($0.1022$ for i.i.d.\ and
$0.1082$ for AR(2) against a realised $0.0957$), and the ratio never
reaches $1.8$ anywhere on the grid, against setting~7's four.

\section{Is the setting-7 tie signal, or degeneracy?}
\label{sec:tie}

A null on a score that has degenerated on most of its cells is not
evidence of no difference. Setting~7 must therefore be checked before its
tie is believed, and the check is available inside the Brier score
itself: compare thresholds whose coverage differs materially.

Moving from $q_{95}$ to $q_{90}$ cuts setting~7's degenerate share from
$60.9\%$ to $41.6\%$ --- the best-covered threshold on the grid. The tie
survives it. The pooled gap is $-0.14 \times 10^{-3}$ ($t\,p = 0.475$),
the dataset count $6/10$, the block count $3/5$; every AR(2) contrast in
\S\ref{sec:tests} is non-significant at $q_{90}$ exactly as at $q_{95}$.
The appendix sweep shows the same across all eleven thresholds: the
setting-7 gap is small and \emph{changes sign} at $q_{91}$, drifting
mildly in i.i.d.'s favour as $p$ rises, with the dataset count decaying
monotonically from $6/10$ to $2/10$. That drift is not evidence for
i.i.d.\ either --- it tracks the over-prediction documented in
\S\ref{sec:coverage}, which penalises AR(2) precisely where realised
exceedance is rarest.

The conclusion is that setting~7's null is a property of the forecasts,
not of the scoring rule. The contrast with setting~5 makes the point
sharply: setting~5 is significant at every threshold from $q_{90}$ to
$q_{98}$ and only loses significance at $q_{99}$ and above, where
$91\%$ of cells are degenerate. Degeneracy does destroy the Brier score's
power --- visibly, at the top of the grid --- but at $q_{90}$ and
$q_{95}$ there is enough left to separate setting~5 comfortably, so the
failure to separate setting~7 is informative.

The control tightens the same argument from the other side, and it does
so on matched coverage. At $q_{90}$ setting~4 leaves $63.3\%$ of cells
degenerate against setting~5's $59.3\%$ --- close enough that the two
nulls cannot be attributed to different amounts of surviving signal ---
yet setting~5 separates at $t\,p = 0.0002$ while setting~4's gap is
$-1.0 \times 10^{-6}$. Nor is the control's flatness confined to one
threshold: the gap stays under $0.08 \times 10^{-3}$ in absolute value
across all eleven, and the count of datasets favouring AR(2) sits at $5$
or $4$ throughout, which is the pattern a genuinely absent effect
produces rather than a suppressed one. Read together, the three settings
order themselves by persistence exactly as the design predicts: nil at
$h^\ast = 6$ days, decisive at $h^\ast = 109$, and, under hyperbolic
decay, a null that \S\ref{sec:coverage} attributes to a level error
rather than to the absence of memory.

\section{Provenance}

Caches: \code{simstudy/output/results/scores4\_hn.RData},
\code{scores5\_hn.RData}, \code{scores7\_hn.RData}, all three with
\code{brier\_threshold\_basis = "full\_series"} and
\code{hn\_prior = TRUE}, written by \code{simstudy/scores.R}. The rule
changed in commit \code{46628f7}; the settings-5 and -7 outputs were
refreshed under it in \code{638f85a}, and setting~4 was fitted and scored
under it in \code{072c07b}.

An untagged \code{scores4.RData} also sits in the same directory. It is
the retired campaign's setting-4 cache --- $N(0, 20)$ prior, two methods,
the block-quantile rule, twenty datasets --- and it is not the file this
report reads. The \code{\_hn} suffix is the only thing separating them on
disk, which is the reason for the tagging convention described next.

The \code{\_hn} suffix is the lambda-prior arm. Since 2026-08-05 every
per-arm artifact carries it --- fits in \code{results\_hn/}, caches in
\code{scores<S>\_hn.RData} --- because the formal experiment runs under
either $\lambda \sim \hn(0, 20)$ (\code{--hn}) or $\lambda \sim N(0, 20)$,
and before tagging the two arms silently overwrote each other. This report
is the HN arm; the N arm has not been fitted.

Tables and their \code{\_body.tex} fragments:
\code{Rscript expA\_breakdown.R --prior=hn "-{}-settings=(4,5,7)"}, which
reads only the three caches and \code{simdata.RData} and writes every
table this report \code{\textbackslash input}s. Each fragment records its arm in a
header comment. No number in this document is transcribed by hand. The script asserts the threshold basis and the absence of NA
cells on load, and cross-checks its $q_{95}$ tests against
\code{expA\_tests\_\{5,7\}.csv}.

Sibling report on the retired rule, including the controlled contrast
between the two rules on these same fits:
\code{tex/timeblock\_brier\_blockq}. Frozen record of the earlier
campaign and of all CRPS results:
\code{tex/timeblock\_expABC\_legacy}.

\appendix
\section{The full threshold grid}

\tbl{expA_qsweep_full}

\noindent Pooled over leads $1$--$15$ at the dataset unit. Setting~5 is
significant on both tests from $q_{90}$ through $q_{98}$ and fails only
at $q_{99}$ and $q_{995}$, where the degenerate share exceeds $91\%$;
the gap decays smoothly with $p$ as the base rate does. Setting~7 never
separates at any threshold. Setting~4 never separates either, and its
gap never leaves the interval $[-0.08, +0.01] \times 10^{-3}$, changing
sign three times inside it; the smallest $t\,p$ on its eleven rows is
$0.44$.


\section{The second-lag contrast across the grid}
\label{sec:sweep-ar2ar1}

The sweep above is AR(2) against i.i.d., which mixes what the first lag buys with what the second adds.  Table~\ref{tab:qsweep-ar2ar1-full} isolates the second lag by sweeping AR(2) against AR(1) over the same grid.  At setting~5 the contrast separates from $q_{90}$ through $q_{98}$, in $10/10$ data sets and $5/5$ blocks throughout, and fails only at $q_{99}$ and $q_{995}$, where more than $91\%$ of cells carry no realised exceedance.  That the second-lag contrast is significant on its own, and over a wider band of the tail than AR(1)$-$i.i.d.\ manages anywhere, is what identifies the second lag rather than temporal pooling in general as the source of the setting-5 result.  Setting~4 never separates ($t\,p$ between $0.15$ and $0.23$); setting~7 never separates either, though its gap is negative at all eleven levels and $7/10$ data sets favour AR(2) over AR(1) at every level up to $q_{97}$ --- a faint and consistent preference that the pooled test cannot resolve at $n = 10$.

\tbl{expA_qsweep_full_ar2ar1}

\end{document}
